\subsection{Autonomous Web Exploration \& Research Agents}
\label{sec:app-research}

Web agents, GUI agents and autonomous research agents constitute three interlinked but distinct trajectories of agentic reasoning systems. Firstly, web agents specialize in navigating online resources, issuing web API calls or browser actions to retrieve dynamic evidence and steer research direction. GUI agents go further by manipulating software interfaces and multi-modal dashboards directly (i.e. clicking, typing, navigating) to execute experiments, data workflows and interface-based tasks. Autonomous research agents sit at the top of this hierarchy, pairing LLM reasoners with scientific workflows, tool-chains and meta-loops to drive hypothesis generation, data synthesis and paper writing. The core connection is a progression of autonomy: first web agents retrieve evidence from online resources, then GUI agents operationalize actions inside software interfaces, and finally autonomous research agents orchestrate full scientific workflows end-to-end.

In this subsection we begin with the foundational layer (Section \ref{sec:app-auto-foundational}), which captures the core capabilities that any autonomous agent must support: perceiving its environment, reasoning about goals, planning actions and grounding those into tool-augmented workflows. Building on these primitives, the self-evolving layer (Section \ref{sec:app-auto-evolving}) examines how agents incorporate feedback, memory and reflection to iteratively refine their behaviors and improve methods over time. Finally, the collective layer (Section \ref{sec:app-auto-multi-agent}) highlights how agents move beyond individual competence into coordination, specialization and emergent collaboration. While web agents, GUI agents and autonomous agents share common themes of goal-directed autonomy, tool-use and iterative improvement, they differ in where they act on, how they manipulate their environment and what goal they aim to achieve.



\subsubsection{Foundational agentic reasoning}


\label{sec:app-auto-foundational}


\paragraph{Planning.}



Planning is essential for web agents because they must decompose long-horizon tasks into manageable steps, adapt to dynamic pages and coordinate tool/invocation strategies. Early work such as WebGPT~\cite{nakano2021webgpt} fine-tuned GPT-3~\cite{brown2020language} to answer open-ended questions via a text-based web-browser interface. Then, various web-based methods deepened the planning paradigm: for example, SEEACT~\cite{zheng2024gpt} explored large multi-modal models as generalists that integrating visual and HTML grounding for web-based tasks, and AutoWebGLM~\cite{lai2024autowebglm} introduced HTML simplification and various learning techniques for open-domain web task decomposition and navigation. These works paved the way for recent systems such as Agent Q~\cite{putta2024agent} that integrate guided MCTS, self-critique and off-policy preference optimization on web-task benchmarks, and set the stage for even more advanced long-horizon web planners such as WebExplorer~\cite{liu2025webexplorer} and WebSailor~\cite{li2025websailor}.




% rl for web agents
In addition, reinforcement learning has become a core tool for improving the decision-making and planning behavior of web-based LLM agents. WebRL~\cite{qi2024webrl} introduces a self-evolving online curriculum that generates new tasks from unsuccessful attempts and trains an outcome-supervised reward model to guide policy optimization. WebAgent-R1~\cite{wei2025webagent} performs end-to-end multi-turn RL, learning web interaction policies directly from online rollouts with binary success rewards. DeepResearcher~\cite{zheng2025deepresearcher} scales RL to real-world web environments, using a multi-agent browsing architecture and exhibiting emergent behaviors such as plan formulation, cross-source corroboration, and self-reflection. Hybrid pipelines like AutoWebGLM~\cite{lai2024autowebglm} combine supervised training with RL fine-tuning to strengthen task decomposition and structured navigation, while Navigating WebAI~\cite{thil2024navigating} combines supervised learning and RL techniques to improve web navigation performance. Methods such as Pangu DeepDiver~\cite{shi2025pangu}, EvolveSearch~\cite{zhang2025evolvesearch}  and WebEvolver~\cite{fang2025webevolver} use RL-based self-improvement, for example by adaptively scaling search depth or jointly training an agent and a world-model-like simulator to improve long-horizon web decision-making. Hierarchical approaches like ArCHer~\cite{zhou2024archer} optimize high-level and low-level policies with a multi-turn hierarchical RL framework, while PAE~\cite{zhou2025proposer} combines a task proposer, an acting agent and an evaluator to support autonomous skill discovery via RL in internet environments. 


% gui agents



% normal planning for gui agents
Planning is a core capability for GUI agents, enabling them to coordinate long, multi-step interactions across applications and operating environments. OS-Copilot~\cite{wu2024copilot} approaches this by treating the desktop as a unified control space in which a generalist agent continually refines its multi-step workflows. Agent S~\cite{agashe2024agent} builds an experience-augmented planning stack that decomposes tasks into sub-goals while retrieving past trajectories and external knowledge to guide action sequencing. InfiGUIAgent~\cite{liu2025infiguiagent} strengthens planning by integrating hierarchical task structuring into a multi-modal backbone, allowing agents to organize GUI procedures at multiple levels of abstraction. MobA~\cite{zhu2025moba} and PC Agent~\cite{liu2025pc} employ hierarchical architectures that separate high-level planning from low-level execution—on mobile and desktop respectively. GUI foundation models such as OS-ATLAS~\cite{wu2024atlas}, OSCAR~\cite{wang2024oscar} and UItron~\cite{zeng2025uitron} further emphasize robust cross-application planning: OS-ATLAS offers a platform-agnostic action model for consistent control, OSCAR maintains state-aware plans that adapt as execution unfolds and UItron unifies offline and online planning within a single general-purpose GUI agent. 


Likewise, reinforcement learning has become a central way to endow GUI agents with planning over long action sequences. End-to-end frameworks such as ARPO~\cite{fanbin2025arpoendtoend} and ComputerRL~\cite{hanyu2025computerrl} directly optimize multi-step GUI trajectories with replay buffers or large-scale online interaction, replacing hand-crafted scripts with learned policies for general desktop control. R1-style and semi-online methods, including UI-R1~\cite{zhengxi2025uir1}, GUI-R1~\cite{run2025guir1}, InfiGUI-R1~\cite{yuhang2025infiguir1} and UI-S1~\cite{zhengxi2025uis1}, start from strong vision–language backbones and then use RL to sharpen action prediction and long-horizon reasoning. A complementary line focuses on where to act by improving visual grounding: GUI-Bee~\cite{yue2025guibee}, SE-GUI~\cite{xinbin2025enhancing}, UIShift/GUI-Shift~\cite{gao2025guishiftenhancingvlmbasedgui} and UI-AGILE~\cite{shuquan2025uiagile} develop RL-based grounding frameworks to help agents reliably localize target elements before executing actions. ZeroGUI~\cite{chenyu2025zerogui} pushes toward fully automated online RL loops, where the agent generates its own tasks and trajectories and improves with zero human annotation, while ComputerRL~\cite{hanyu2025computerrl} scales end-to-end online RL in large distributed desktop environments. AgentCPM-GUI~\cite{zhong2025agentcpmgui} couples supervised pre-training with reinforcement fine-tuning to strengthen decision quality on mobile apps. Finally, foundation-style GUI agents such as AutoGLM~\cite{xiao2024autoglm} and Mobile-Agent-v3~\cite{jiabo2025mobileagentv3} serve as general backbones that unify perception, grounding and action, and are trained or fine-tuned with scalable RL frameworks to align long-horizon GUI planning with real-world success signals.

% research agents
For autonomous research agents, planning modules translate abstract goals into actionable research itineraries.  For example, Agent Laboratory~\cite{schmidgall2025agentlaboratoryusingllm} organizes work into three structured stages, namely literature review, experimentation and report writing, and supports the workflow with tool-hooks that automate code execution, experiment runs and documentation. GPT Researcher~\cite{Elovic2023gptresearcher} uses a plan → research → write cycle, where a dedicated planner drafts the outline, retrieval/analysis agents gather evidence and a writer compiles the final report. Chain of Ideas~\cite{li2024chainideasrevolutionizingresearch} retrieves literature into a chain structure to reflect domain progression and support ideation via experiment design, whereas IRIS~\cite{garikaparthi2025iris} performs hypothesis exploration via Monte Carlo Tree Search to expand promising branches before committing to downstream tasks. Broader variants include ARIA~\cite{ramirezmedina2025acceleratingscientificresearchmultillm} and NovelSeek~\cite{novelseekteam2025novelseekagentscientist} that automate the research workflow with a complete literature search, hypothesis generation and experiment planning cycle.




\paragraph{Tool-Use.}



% web agents
For web agents, tool-use abilities underpins execute plans in realistic, dynamic environments. For example, WebVoyager~\cite{he2024webvoyager} systematizes multi-modal execution by building an end-to-end agent that operates on real websites. On the interaction side, BrowserAgent~\cite{zhang2025browseragent} makes the action space more human-like, defining a compact set of browser primitives (e.g., click, scroll, type) and coupling them with an explicit memory mechanism to maintain key conclusions across steps, yielding strong gains on multi-hop QA benchmarks. Finally, methods like WALT~\cite{prabhu2025walt} and pipeline-oriented systems such as WebDancer~\cite{wu2025webdancer} and WebShaper~\cite{tao2025webshaper} push tool use from mere execution toward tool discovery and data-centric interaction. Specifically, WALT teaches agents to reverse-engineer reusable tools from website functionality, while WebDancer and WebShaper embed web actions inside multi-turn information-seeking and dataset-synthesis loops, respectively.


% gui agent tool-use
Tool use is another core capability for GUI agents, enabling them to invoke system functions and application features as structured tools. As pioneering systems, AutoDroid~\cite{wen2024autodroid} automatically analyzes Android apps to construct functionality-aware UI abstractions that LLM agents can reason over as capabilities rather than raw layouts, while its successor AutoDroid-V2~\cite{wen2025autodroid} re-frames mobile UI automation as LLM-driven code generation, with an on-device small language model emitting executable scripts for a local interpreter. MobileExperts~\cite{zhang2024mobileexperts} models each expert as a tool-capable specialist and uses a dual-layer controller to select which expert and its associated tool-set to invoke at different stages of a mobile workflow. AgentStore~\cite{jia2025agentstore} pushes this idea to the platform level by treating heterogeneous agents themselves as tools: a MetaAgent uses AgentTokens to route operating-system subtasks to the most suitable specialized “tool-agent” through a unified interface. OS-Copilot~\cite{wu2024copilot} and OSCAR~\cite{wang2024oscar} integrate rich system-level tools into unified computer-control frameworks, so that complex desktop tasks are expressed as sequences of tool calls. OS-ATLAS~\cite{wu2024atlas} complements these systems with a foundation action model that offers robust cross-platform GUI grounding, serving as a reliable actuator layer for downstream tool-using agents. Finally, SeeClick~\cite{cheng2024seeclick} strengthens the execution stack by pre-training a visual GUI agent for GUI grounding, improving the ability to locate the correct on-screen elements from instructions.


Specialized tools can expand an autonomous research agent’s capabilities beyond pure text, allowing more fine-grained ability. For instance, Agentic Reasoning~\cite{wu2025agentic} automatically routes queries to appropriate tool modules like code execution, web search and structured memory agents when the main LLM detects a gap in reasoning; while Webthinker~\cite{li2025webthinkerempoweringlargereasoning} empowers autonomous web exploration and page navigation during long-horizon investigations, by interleaving reasoning, search and draft-writing with a web explorer module. PaperQA~\cite{lála2023paperqaretrievalaugmentedgenerativeagent} and its follow-up synthesis agent~\cite{skarlinski2024languageagentsachievesuperhuman} integrate PDF parsing and citation-level grounding to produce verifiable answers and literature syntheses, while Scideator~\cite{radensky2025scideatorhumanllmscientificidea} provides an IDE-style tool-chain that combines paper facets with novelty checks for real-time brainstorming. In addition, DeepResearcher~\cite{zheng2025deepresearcher} shows that reinforcement learning over real-web interactions improves deep-research efficiency and quality, with emergent behaviors such as plan refinement and cross-source corroboration. 


Execution components ground high-level reasoning in code, simulations or laboratory protocols to produce verifiable scientific outcomes. Agent Laboratory~\cite{schmidgall2025agentlaboratoryusingllm} executes experiments specified in declarative configuration files by orchestrating external toolchains, while Agentic Reasoning~\cite{wu2025agentic} integrates a coding agent that executes Python alongside web search and structured memory, feeding the results back into the reasoning process. MLR-Copilot~\cite{li2024mlrcopilotautonomousmachinelearning} turns research plans into runnable implementations via an ExperimentAgent that leverages retrieved prototype code, runs experiments, and iteratively debugs implementations. Dolphin~\cite{yuan2025dolphin} closes the loop by generating ideas, implementing them through code templates with traceback-guided debugging, executing experiments, and using the analyzed results to steer the next research cycle. The AI Scientist~\cite{chris2024the} automates end-to-end ML experiments, i.e. generating ideas, writing code, executing experiments and visualizing results, so that observed outcomes can guide subsequent runs, while The AI Scientist-v2~\cite{yamada2025ai} adds a dedicated experiment manager and progressive agentic tree search to prioritize and schedule experiment branches. Most recently, NovelSeek~\cite{novelseekteam2025novelseekagentscientist} introduces a unified closed-loop multi-agent framework that spans hypothesis generation, idea-to-methodology construction, and multi-round automated experiment execution with feedback across diverse scientific domains.


\paragraph{Search and Retrieval.}


Search and retrieval lie at the heart of what differentiates web agents from static language models: they must locate, synthesize and refactor web-scale information in dynamic environments. WebExplorer~\cite{liu2025webexplorer} tackles this by generating challenging information-seeking trajectories and training agents to interleave a search tool and a browse tool over many turns, resulting in improved multi-step retrieval policies on complex benchmarks. WebSailor~\cite{li2025websailor} likewise focuses on information-seeking under extreme uncertainty, constructing high-uncertainty search tasks and using a two-stage post-training pipeline  to instill uncertainty-reducing search strategies for long-horizon web tasks. INFOGENT~\cite{reddy2025infogent} also performs  multi-query search across diverse web sources, enabling comprehensive information retrieval beyond task completion. For retrieval-augmented generation applications, RaDA~\cite{kim2024rada} explicitly disentangles web-agent planning into Retrieval-augmented Task Decomposition and Retrieval-augmented Action Generation, so that each high-level subgoal and concrete action is conditioned on fresh search results while respecting context limits. In addition, GeAR~\cite{shen2025gear} advances retrieval itself by augmenting a base retriever with graph expansion and an agent framework, enabling multi-hop passage retrieval along graph-structured evidence chains. Finally, WebRAGent~\cite{anonymous2025webragent} exemplifies retrieval-augmented generation for web agents by retrieving past trajectories and external knowledge into a multi-modal RAG policy.



% gui agents
Several GUI agents use retrieval capability to inject external experience or knowledge at inference time. Synapse~\cite{zheng2023synapse} maintains an exemplar memory of abstracted trajectories and, for each new task, retrieves similar past trajectories as in-context plans, substantially improving multi-step decision-making. LearnAct~\cite{liu2025learnact} builds a three-agent pipeline that mines human demonstrations into a knowledge store and retrieves the most relevant instructions to guide mobile GUI execution on unseen and diverse tasks. MobileGPT’s Explore–Select–Derive–Recall framework~\cite{lee2023explore} equips a phone agent with human-like app memory, storing modular procedures that can be recalled and recomposed when similar tasks reappear. TongUI~\cite{zhang2025tonguiinternetscaletrajectoriesmultimodal} turns large-scale multi-modal web tutorials into the GUI-Net trajectory corpus, effectively giving agents a large offline memory of how humans operate hundreds of apps across multiple operating systems. RAG-GUI~\cite{xu2025retrieval} makes retrieval explicit at inference time by querying web tutorials and generating textual guidelines that are fed into any VLM-based GUI agent as step-by-step hints. WebRAGent~\cite{anonymous2025webragent} shows a related pattern in web automation, combining a multi-modal retriever with a web agent so that each action is conditioned on retrieved guidance.


% research agents
Search modules probe the research landscape to surface relevant papers and passages, enriching context and grounding subsequent reasoning. WebThinker~\cite{li2025webthinkerempoweringlargereasoning} equips large reasoning models with a Deep Web Explorer module for autonomous web search and page navigation, and uses an Autonomous Think–Search–and–Draft strategy with RL-based training to decide when to browse and what to extract during long-horizon tasks. DeepResearcher~\cite{zheng2025deepresearcher} scales end-to-end training on the real web via reinforcement learning in live search environments, optimizing the iterative think–search loop and exhibiting emergent behaviors such as plan formulation, cross-source corroboration, and self-correction over multi-step research trajectories. Retrieval-centric agents like PaperQA~\cite{lála2023paperqaretrievalaugmentedgenerativeagent} and its successor PaperQA2~\cite{skarlinski2024languageagentsachievesuperhuman} demonstrate that tightly coupling full-text retrieval with generation can substantially improve scientific QA accuracy while preserving cited provenance for literature synthesis.

In research settings, retrieval-augmented generation (RAG) grounds idea generation and analysis in freshly retrieved and citable passages. For example, GPT Researcher~\cite{Elovic2023gptresearcher} is an autonomous research-agent that retrieves sources and generates reports with citations, enabling traceability of claims to evidence. Chain of Ideas~\cite{li2024chainideasrevolutionizingresearch} organizes relevant literature into a chain-structured scaffold that mirrors a field’s progressive development, thereby guiding retrieval and ideation toward subsequent links in the argument. Meanwhile, Scideator~\cite{radensky2025scideatorhumanllmscientificidea} extracts key paper facets (e.g., purpose, mechanism, evaluation) and leverages them to drive targeted retrieval and recombination of ideas for identifying methodological or evidentiary gaps.




\subsubsection{Self-evolving agentic reasoning}
\label{sec:app-auto-evolving}



Effective self-evolving abilities enable these autonomous agents to adapt their behavior over time, retain crucial task context across interaction cycles and incrementally refine planning and execution strategies. The following paragraphs review how memory, feedback and self-reflection mechanisms support this continual improvement across these agent families, turning interaction from a one-shot pipeline into an iterative learning loop.





% web agents
\paragraph{Memory.}


Memory modules transform brittle, single-pass web interactions into reusable experience. For example, Agent Workflow Memory (AWM)~\cite{Wang2024AgentWM} induces reusable workflows from successful trajectories and retrieves them to guide future tasks, while ICAL~\cite{sarch2024vlm} distills noisy trajectories into high-level verbal and visual abstractions that are stored as a memory of multimodal experience and later injected into prompts. Control-oriented designs such as BrowserAgent~\cite{zhang2025browseragent} maintain explicit histories of past actions and intermediate conclusions in the agent’s context, instead of only re-encoding the current page view. GLM-based agents like AutoWebGLM~\cite{lai2024autowebglm} and AgentOccam~\cite{yang2024agentoccam} emphasize compressed page representations, using HTML simplification and carefully tuned observation spaces so that the agent’s prompt contains a shorter, more informative view of the state, with past steps preserved through the usual action–observation history. More integrated frameworks like LiteWebAgent~\cite{zhang2025litewebagent} expose planning, memory and tree search as modular components, and can plug in workflow memories together with search traces for long-horizon reuse.



% gui agents
Recent GUI agents adopt explicit memory modules that store and retrieve task-relevant information during long-horizon execution. Earlier work such as MobileGPT~\cite{lee2024mobilegpt} equips a mobile assistant with human-like app memory: it decomposes procedures into modular sub-tasks that are explored, selected, derived, and then stored so they can be recalled and reused instead of being re-discovered from scratch. Chain-of-Memory (CoM)~\cite{gao2025chain} incorporates short- and long-term memory by recording action descriptions and task-relevant screen information in a dedicated memory module, enabling cross-application navigation to track task state. More recent systems build increasingly structured memories: MobA’s multifaceted memory module~\cite{zhu2025moba} maintains environment- and user-level traces that an adaptive planner retrieves when refining mobile task plans, while MGA~\cite{cheng2025mga} represents each step as a triad of current screenshot, spatial layout, and a dynamically updated structured memory that summarizes past transitions, mitigating error accumulation in long chains of actions. Mobile-Agent-E~\cite{wang2025mobileagente} adds a persistent long-term store of tips and shortcuts distilled from prior trajectories, so later plans can call reusable guidance and subroutines instead of relearning them. Mirage-1~\cite{xie2025mirage} similarly organizes experience into a hierarchical skill memory that a planner can retrieve as reusable building blocks for new GUI tasks. 



% research agents
Long-term memory is crucial for autonomous research agents because it enables accumulation and reuse of prior knowledge, fostering continuity across research cycles. For example, Agent Laboratory~\cite{schmidgall2025agentlaboratoryusingllm} retains prior experiment code, results, and interpretation across its multi-phase workflow, enabling later stages to build on earlier work. GPT Researcher~\cite{Elovic2023gptresearcher} generates reports with embedded citations and provides context for planning and extension of research topics. Chain of Ideas~\cite{li2024chainideasrevolutionizingresearch} structures relevant literature into a chain scaffold that reflects a field’s progression and can be revisited as new evidence arises. The AI Scientist-v2~\cite{yamada2025ai} incorporates a progressive agentic tree-search approach that enables branching, backtracking and follow-up experimentation across iterations.




\paragraph{Agentic Feedback and Reflection.}


Modern web agents treat interaction as a continual learning process, using feedback signals and reflection modules to refine their reasoning and recover from failures over time. Agent Q~\cite{putta2024agent} combines guided Monte Carlo tree search with a self-critique stage, so that rollouts provide not only action sequences but also preference-style supervision. \textsc{ReAP}~\cite{azam2025reflection} makes reflection explicit by treating it as a retrieval problem: it stores task–reflection key–value pairs summarizing what was learned from past trajectories, then, at inference time, retrieves the most relevant reflections and appends them to the agent’s prompt to guide planning on new web-navigation tasks. Agent-E~\cite{azam2024multimodal} introduces an automatic validation pipeline that detects execution errors across text and vision, and then triggers self-refinement, enabling agents to iteratively correct their own workflows. Recon-Act~\cite{he2025recon} uses a dual-team architecture in which a Reconnaissance team extracts generalized tools from successful and failed trajectories, and an Action team applies these tools to re-plan tasks, forming a closed feedback loop. INFOGENT~\cite{reddy2025infogent}, on the other hand, leverages aggregator feedback to iteratively refine navigation and search strategies based on identified information gaps. And WINELL~\cite{reddy2025winell}, as a updating web agent, relies on feedback from the aggregation process to adapt subsequent searches and update selection during continuous operation. Finally, self-reflective search agents such as WebSeer~\cite{he2025webseer} integrate explicit self-reflection signals into reinforcement learning, constructing reflection-annotated trajectories and a two-stage training framework so that mis-solved or uncertain cases become targeted feedback that deepens future search and reasoning.

GUI agents also integrate explicit reflection so they can critique and repair their own plans. Early computer-control systems with structured reflection, for example, a zero-shot desktop control agent with structured self-reflection loops~\cite{li2023zero}, provides conceptual templates that later GUI agents adapt to visual, multi-application settings. GUI-Reflection~\cite{wu2025gui} instantiates this idea end-to-end: it builds a reflection-oriented task suite, automatically synthesizes error scenarios from existing successful trajectories, and adds an online reflection-tuning stage so multi-modal GUI models learn to detect failures, reason about causes, and generate corrective actions without human annotation. History-Aware Reasoning (HAR)~\cite{wang2025history} treats long-horizon GUI automation as a reflective learning problem, constructing reflective learning scenarios, synthesizing tailored correction guidelines, and designing a hybrid RL reward so the agent acquires episodic reasoning knowledge from its own errors and shifts from history-agnostic to history-aware reasoning. MobileUse~\cite{li2025mobileuse} introduces hierarchical reflection on mobile devices, where the agent self-monitors at the action, subtask, and task level and triggers reflection on demand, pausing only when needed to diagnose and recover. InfiGUIAgent~\cite{liu2025infiguiagent} integrates hierarchical and expectation–reflection reasoning in a second training stage, enabling the agent to run expectation–reflection cycles that compare expected and actual outcomes and revise multi-step plans when they diverge. Mobile-Agent-E~\cite{wang2025mobileagente} embeds an Action Reflector and Notetaker that evaluate executed steps and write refined Tips and Shortcuts back into persistent long-term memory, forming a self-evolution loop where the agent’s behavior is progressively refined from accumulated experience. 


For autonomous research agents, learning from outcomes is essential to improve reasoning and experimental reliability over time. CycleResearcher~\cite{weng2025cycleresearcherimprovingautomatedresearch} couples a research agent with a reviewer agent that provides automated peer-review feedback, and uses an iterative preference-training loop so the research agent can refine future drafts and decisions. MLR-Copilot~\cite{li2024mlrcopilotautonomousmachinelearning} monitors execution results and human comments during experiment implementation and execution, using these signals to iteratively refine code, configurations and even upstream hypotheses. Dolphin~\cite{yuan2025dolphin} implements a closed-loop auto-research framework in which generated code is run on benchmarks and exception-guided debugging plus outcome analysis feed back into idea generation and implementation, pruning unproductive paths. At the search–reasoning interface, DeepResearcher~\cite{zheng2025deepresearcher} optimizes query, browsing, and answering policies via reinforcement learning on real-web trajectories, with outcome rewards inducing behaviors such as planning, cross-validation, and self-reflection. Agentic Deep Research~\cite{weizhi2025from} further emphasizes reward design for reasoning-driven search, arguing that principled incentives over answer quality and reasoning traces provide structured signals that improve downstream synthesis in deep-research agents.




\subsubsection{Collective multi-agent reasoning}
\label{sec:app-auto-multi-agent}


% web agent
Collective multi-agent reasoning for web agents reframes browser use as cooperation among specialized roles rather than a single monolithic policy. WebPilot~\cite{zhang2025webpilot} models web task execution as a multi-agent system with a global planning agent that decomposes tasks and local MCTS-based executors that solve subtasks, jointly steering search in complex web environments. INFOGENT~\cite{reddy2025infogent} organizes web information aggregation into a Navigator, Extractor, and Aggregator, so exploration, evidence extraction, and synthesis are handled by distinct cooperating agents with feedback from the Aggregator to guide future navigation; WINELL~\cite{reddy2025winell} leverages agentic web search to plan and execute iterative information gathering for discovering timely factual updates relevant to a target Wikipedia article.. Recon-Act~\cite{he2025recon} adopts a Reconnaissance–Action paradigm in which a Recon team analyzes successful and failed trajectories to derive generalized tools or hints, and an Action team re-plans and executes with this evolving toolset. PAE~\cite{zhou2025proposer} uses three roles, namely a task proposer, an acting web agent and a VLM-based evaluator, to autonomously generate vision-based web tasks and feed success signals back into the policy via RL. Hierarchical web agents such as Agent-E~\cite{abuelsaad2024agent} and Plan-and-Act~\cite{wang2023plan} similarly separate a high-level planner from a browser-navigation agent, enabling structured plan–execute cooperation. At a more conceptual level, Agentic Web~\cite{yang2025agentic} envisions the internet as an agentic web of interacting agents and analyzes how coordination, communication protocols and economic incentives shape such ecosystems, while Agentic Deep Research~\cite{weizhi2025from} frames information seeking as iterative feedback loops of reasoning, retrieval, and synthesis that can be instantiated by single- or multi-agent web research systems.



% gui agent
Multi-agent designs for GUI agents typically decompose “using a computer” into cooperating roles that plan, perceive, decide and execute. COLA~\cite{zhao2025cola} instantiates a scenario-aware task scheduler, a planner, a decision-agent pool, an executor, and a reviewer, so UI tasks are split into basic capability units and routed to domain-specialized agents rather than a single monolith. On mobile, Mobile-Agent-v2~\cite{wang2024mobile} adopts a tri-role pattern with planner, decision, and reflection agents for progress navigation, local action selection, and error correction, while Mobile-Agent-E~\cite{wang2025mobileagente} further builds a hierarchical stack with a Manager and four subordinate agents (i.e. Perceptor, Operator, Action Reflector, Notetaker) plus a self-evolution module that learns long-term Tips and Shortcuts from experience. Mobile-Agent-V~\cite{wang2025mobileagentvvideoguidedapproacheffortless} similarly employs a video agent, decision agent, and reflection agent to coordinate multi-modal perception and execution, and MobileExperts~\cite{zhang2024mobileexperts} dynamically forms teams of expert agents with a dual-layer planner that allocates subtasks to tool-specialized experts. SWIRL~\cite{lu2025swirl} makes this structure explicit for RL, training a Navigator that converts language and screen context into structured plans and an Interactor that grounds those plans into atomic GUI actions within a multi-agent RL workflow. PC Agent~\cite{liu2025pc} uses separate planning and grounding agents in a two-stage pipeline for desktop automation, illustrating how multi-agent decomposition can improve long-horizon PC control. 

% research agent
To facilitate autonomous research agents, multi-agent collaboration enables a single model’s linear workflow to become a coordinated research group: specialized agents operate in parallel, exchange intermediate artifacts through explicit interfaces, and provide adversarial or complementary feedback to improve both creativity and rigor. For example, AgentRxiv~\cite{schmidgall2025agentrxivcollaborativeautonomousresearch} coordinates author, reviewer, and editor agents that iteratively refine manuscripts and share evolving artifacts across virtual “labs.” ARIA~\cite{ramirezmedina2025acceleratingscientificresearchmultillm} instantiates a role-structured multi-LLM team that searches, filters, and synthesizes scientific literature into actionable experimental procedures. Earlier multi-agent designs such as CAMEL~\cite{qi2023largelanguagemodelszero} demonstrate how cooperative role-play with tool access can enhance hypothesis generation and task decomposition. In experimental sciences, Coscientist~\cite{boiko2023autonomous} integrates planning, robotic instrument control, and analysis into a multi-agent closed loop that autonomously designs and executes wet-lab experiments. Finally, TAIS~\cite{liu2024toward} defines a hierarchical team, namely project manager, data engineer and domain expert, that jointly discovers disease-predictive genes from expression data through coordinated division of labor.
